% 本文件由 LaTeX 排版 Agent 根据有效 translation.json 自动生成。
% author=Tertullian; work_id=0305; title=致斯卡普拉

\chapter{致斯卡普拉}

% structure: p0001 source=h1 target=omitted reason=page-title-already-rendered-by-parent

\EditorialNote{斯卡普拉是迦太基的总督，虽然其年代是推测的（公元217年），但这部著作提供了其时代和环境的宝贵指标。它写于塞维鲁死后，第四章中提及了他；在拜占庭（公元196年）毁灭之后，第三章中提到了这一点；阿利克斯博士推测，在尤提卡的黑暗日（公元210年）之后，他认为同一章中提到了这一点。第四章中提到的辛丘斯·塞维鲁于公元198年被塞维鲁处死。}\par

% structure: p0003 source=h2 target=section reason=relative-heading-rank; base=h2

\section{第一章}

我们并非因遭受来自人们无知的迫害而极度不安或惊慌；因为我们加入此教派时，已完全接受其盟约的条款，以致作为生命不属于自己之人，我们投入这些争战，我们的渴望是获得天主应许的赏报，我们的恐惧是祂以非基督徒的生活相威胁的灾祸会临到我们。因此，我们毫不退缩地与你们最大的愤怒搏斗，甚至自愿前来应战；定罪比开释更令我们喜悦。因此，我们将此小册子寄给你们，并非为自己忧虑，而是非常关心你们和所有我们的仇人，更不用说我们的朋友。因为我们的宗教命令我们甚至爱我们的仇人，并为那些迫害我们的人祈祷，追求其特有的成全，并在其门徒中寻求一种高于世上普通善良的更高类型。因为众人都爱那爱他们的人；唯独基督徒拥有爱那恨他们之人的特质。因此，为你们的无知而哀恸，怜悯人类的错误，并展望那每一天都显出威胁迹象的未来，我们也有必要以这种方式出来，以便将你们不愿公开聆听的真理摆在你们面前。\par

% structure: p0005 source=h2 target=section reason=relative-heading-rank; base=h2

\section{第二章}

我们崇拜独一的天主，关于祂的存在和属性，自然教导所有人；在祂的闪电和雷轰下你们战栗，祂的恩惠促成你们的幸福。你们认为其他也是神，而我们晓得它们是魔鬼。然而，按照自己的信念崇拜，是一项基本人权，是自然的特权：一个人的宗教既不伤害也不帮助另一个人。宗教决不包括强迫宗教——自由意志而非强迫应引导我们——甚至祭品也需要出自自愿之心。你们强迫我们献祭，对你们的神并没有真正的服务。因为他们不能渴望来自不情愿者的祭品，除非他们被争斗的精神所驱动，这完全是非神的。因此，真天主将祂的祝福同等地赐给恶人和祂的选民；为此，祂设立了永久的审判，那时感恩的和不感恩的必须在祂的审判台前站立。然而，你们从未发现我们——虽然你们认为我们是亵圣的恶徒——有任何偷窃，更不用说任何亵圣了。但你们圣殿的强盗，他们都指着你们的神起誓并崇拜他们；他们不是基督徒，然而正是他们被发现犯有亵圣行为。我们没有时间展开讲述你们的神被他们自己的崇拜者以多少其他方式嘲弄和蔑视。同样，叛国罪也错误地加在我们的指控上，尽管从未有人能在基督徒中找到\Person{阿尔比努斯}、\Person{尼格尔}或\Person{卡西乌斯}的追随者；而那些曾指着皇帝的守护神起誓、为他们的安全献祭并许愿、经常宣布对基督门徒的定罪的人，直到今日仍被发现是帝国王座的叛徒。基督徒与任何人都无仇，更不对\Place{罗马}皇帝有仇，他知道他是被他的天主所命定，因此不能不敬爱他；并且他必须渴望他的福祉，以及他所统治的帝国的福祉，只要世界尚存——\Place{罗马}将持续多久。因此，我们向皇帝致以我们合法且对他有益的尊敬；视他为仅次于天主的人，他从天主领受了他一切权能，仅低于天主。这将符合他自己的愿望。因为这样——既然仅低于真天主——他比所有其他都伟大。因此他比诸神本身更伟大，就连他们也服从他。因此，我们为皇帝的安全献祭，但是献祭给我们的天主和他的天主，并按照天主所吩咐的方式，以单纯的祈祷。因为宇宙的创造主不需要香气或血。这些东西是魔鬼的食物。但我们不仅拒绝这些恶灵：我们战胜他们；我们每天使他们受轻蔑；我们从受害者身上驱逐他们，正如众多人可以作证。所以我们更加为皇帝的福祉祈祷，如同那些向能够赐予它的那一位寻求的人。并且人们会认为，你们应当十分清楚，我们据以行动的宗教体系所灌输的是天主的忍耐；因为，尽管我们人数如此众多——几乎在每个城市构成多数——我们却如此安静谦逊地行事；我或许可以说，我们作为个人而非有组织的团体为人所知，仅因我们从前恶习的改正而引人注目。因为我们绝不会对加在我们身上的我们自己所希望的事感到不满，或以任何方式图谋亲手的报复，我们等待从天主而来的报复。\par

% structure: p0007 source=h2 target=section reason=relative-heading-rank; base=h2

\section{第三章}

然而，正如我们已指出的，令我们感到痛苦的是，没有一个邦国可以免受流基督徒之血的罪孽惩罚；正如你确实看到，在\Person{希拉里安}任总督时所发生的事：当我们为死者安葬地点有所争议时，有人喊说：\Quote{没有\Emph{\OriginalTerm{场地}{areæ}}——基督徒没有墓地}，结果他们自己的\Emph{\OriginalTerm{场地}{areæ}}，即他们的打谷场，竟荒废了，因为他们没有收成。至于去年的雨水，它们意在提醒人们什么，那是显而易见的——无疑是指古代淹没人类不信与邪恶的洪水；至于近日彻夜悬于\Place{迦太基}城墙上空的火焰，亲眼看见的人知道它们预示什么；而此前雷声的轰鸣，被它们硬了心肠的\Emph{那些人}可以述说。这一切都是天主即将发怒的征兆，我们必须以各种方式加以宣扬公布；同时我们必须祈祷这愤怒仅限于局部。那些对这些预兆作其他解释的人，\Emph{他们}必有一日要经历其普遍而最终的形态。同样，在\Place{乌提卡}首府，太阳几乎熄灭，这是一个预兆，不可能由平常的日食引起，因为太阳正值其高度和宫位。你们有占星家，可以咨询他们。我们还可以指出一些行省总督的死亡，他们在临终时痛苦地回忆起迫害基督追随者的罪孽。\Person{维吉利乌·撒杜尔尼诺}，他最先在此地用刀剑对付我们，失明了。\Place{卡帕多西亚}的\Person{克劳狄乌·路济乌·赫尔米尼亚努斯}，因他的妻子成了基督徒而大怒，曾极其残忍地对待基督徒：后来他独自留在宫中，患了传染病，浑身溃烂生出活蛆，有人听见他喊道：\Quote{不要让人知道，免得基督徒欢喜，基督徒的妻子们受到鼓励。}后来，他认识到自己的错误，即用施加的酷刑诱使许多人放弃坚定，而他本人几乎像一个基督徒死去。在降临\Place{拜占庭}的灾祸中，\Person{凯基利乌·卡佩拉}不禁喊道：\Quote{基督徒，欢喜吧！}是的，那些自以为可以免受惩罚的迫害者，必逃不过审判之日。至于你，我们真诚希望这仅是一个警告：当你在\Place{阿德鲁梅图姆}将\Person{马维鲁斯}判决喂野兽之后，你便立即遭受了那些麻烦，而如今同样为此原因，你被要求进行血算。但不要忘记未来。\par

% structure: p0009 source=h2 target=section reason=relative-heading-rank; base=h2

\section{第四章}

我们自身无所畏惧，并非要恐吓你，而是若可能，愿藉警告众人不要与天主争斗，来拯救所有的人。你尽可履行你职务的职责，但仍当记念人性的要求；即使只基于你自身也难逃惩罚这一理由，你本该如此。因为你的职权无非是定那些认罪者之罪，并将否认者交付酷刑，难道不是吗？你看，你如何违逆你的指令，强迫认罪者否认。事实上，我们承认时你拒绝立即定罪，正是对我们无辜的承认。你若竭力要根除我们，那你所攻击的正是无辜。然而，有多少比你更果决更残忍的统治者，曾设法完全摆脱这类案件——如\Person{Cincius Severus}，他本人曾在\Place{Thysdris}提出补救之策，指出基督徒应如何回答方可获释；如\Person{Vespronius Candidus}，他释放了一位基督徒，理由是满足同胞会破坏社群的和平；如\Person{Asper}，他在一个因轻微酷刑而背弃信仰者的案件中，并未强迫献祭，并曾在法庭的辩护人和顾问面前承认，他对不得不处理这样的案件感到懊恼。又如\Person{Pudens}，他立刻释放了被带到面前的一位基督徒，从起诉书中看出这是一桩恶意的指控；他撕碎文件，拒绝在没有控告者在场时听审，因为这不符合皇帝的命令。这一切本可正式由那些辩护人提交给你，他们自身也对我们负有义务，尽管在法庭上他们随己意发言。他们中有一位书记官，曾因恶灵附体而跌倒，被从痛苦中释放；另一位亲戚，以及第三位的小儿子也是如此。有多少显贵之人（更不必说平民）曾被从魔鬼手中解救、疾病得愈！甚至\Person{Severus}本人，\Person{Antonine}的父亲，也曾恩慈地记念基督徒；他寻找基督徒\Person{Proculus}，别号\OriginalTerm{Torpacion}{Torpacion}，\Person{Euhodias}的管家，因他曾用敷油治愈过\Person{Severus}，便将他留在宫中直到去世。\Person{Antonine}自幼受基督徒的养育，也与这人熟识。\Person{Severus}深知为基督徒的最高贵男女，不仅允许他们安然无恙，还公开为他们作有利的证词，并将他们从暴怒的民众手中归还给我们。\Person{Marcus Aurelius}在远征日耳曼时，也因他军中基督徒士兵向天主的祈祷，在那众所周知的干渴中获得了雨水。的确，何时我们的跪拜和禁食未能驱散干旱？在这样的时刻，人民呼求\Quote{诸神之神，唯一全能者}，以\OriginalTerm{Jupiter}{Jupiter}之名，为我们的天主作了见证。此外，我们从不否认托付给我们的寄托；我们从不玷污婚姻之床；我们忠诚地对待受我们监护的人；我们帮助贫困者；我们不以恶报恶。至于那些虚假声称属于我们、我们也拒绝承认的人，让他们自己回答吧。简言之，还有谁在其他方面控诉我们？基督徒除了他自身团体的事务，还致力于什么？这团体在其整个悠久的存在中，从未有人证实犯有被指控的乱伦或残暴之罪。正是因这非凡的无罪、如此伟大的正直、公义、纯洁、忠诚、真理，以及生活的天主，我们被投入火焰；因为这种惩罚，你们通常既不施加于亵圣者，也不施于确凿的公敌，也不施于你们那许多叛国者。甚至如今，我们的人在\Place{Legio}和\Place{Mauritania}的总督手下遭受迫害；但只是用刀剑，正如最初所命定我们应受苦一样。然而，我们的冲突越大，我们的赏报也越大。\par

% structure: p0011 source=h2 target=section reason=relative-heading-rank; base=h2

\section{第五章}

你们的残忍是我们的荣耀。你只管看，在忍受这样的事时，我们是否感到被迫冲出去战斗，只为证明我们并不畏惧它们，反而甚至邀请它们施加于我们。当\Person{Arrius Antoninus}在\Place{亚细亚}严厉行事时，全省的基督徒团结一致，出现在他的审判座前；他下令将少数人拉出去处决，对余下的说：\Quote{可怜的人哪，你们若想死，有的是悬崖或绞索。}我们若在此地同样行事，当成千上万的人——如此众多的男女、每一种性别、年龄和阶层——出现在你面前时，你将如何处置？需要多少火、多少剑？\Place{迦太基}本身将何等痛苦，你将不得不屠戮它十分之一的人口，因每个人在那里认出自己的亲戚和同伴，他看到那里或许有你同一等级的人、贵妇和城中所有首领人物，以及你圈内之人的亲属或朋友？请宽恕你自己，若不宽恕我们可怜的基督徒！请宽恕\Place{迦太基}，若不宽恕你自己！请宽恕这行省，你的意图已使之同时遭受士兵和私人仇敌的威胁与勒索。\par

\Emph{我们}除天主外别无主人。\Emph{祂在}你面前，你无法隐藏，但你对祂不能造成任何伤害。而你们视为主人的，不过是人，有一天他们自己也要死去。然而这团体永不止息，因为你要确信，正是在它看似倾覆之时，它被建立成更大的力量。所有目睹其殉道者高尚忍耐的人，如同受震撼而疑虑，便燃起探究此事的渴望；一旦认识真理，便立刻成为它的门徒。\par

\SourceRecord{Translated by S. Thelwall.；Ante-Nicene Fathers；Vol. 3.；Edited by Alexander Roberts, James Donaldson, and A. Cleveland Coxe.；Buffalo, NY: Christian Literature Publishing Co.,；1885.；<http://www.newadvent.org/fathers/0305.htm>.}
